\scanpage{24}
\[
\unsure{\mathbb P|P}^1\times\mathbb P^1\simeq Q^2\subset Q^1\simeq P^1,\qquad
\frac{n-1}{2}+1=\frac{n+1}{2}
\]
\diagram{\begin{tikzcd}[ampersand replacement=\&]
\cdot\arrow[r,dash] \& H^0(Q^1)\arrow[r] \& H^2(\\
{} \& \unsure{\checkmark} \& {}
\end{tikzcd}}
\[
\begin{aligned}
&& n=2\nu\\
H^2(Q^n)&\mathrel{\rule{1.6em}{.35pt}}& H^{2\nu-2}(Q^n)
\end{aligned}
\]
\[
\begin{aligned}
\xi^{0}&\ \mathord{-}\ \xi^{\nu-1}&\xi^\nu&\quad \deleted{\unsure{\xi}}\xi^{\nu+1\,\unsure{[\tfrac12|[2|/2|[\xi}}\ \cdots\ \xi^{2\nu}\\
&&\unsure{l|h|\ell}&\\
\xi^{\unsure{1|4}}_{2}&\quad \deleted{\cdots\ \xi_{\nu}\qquad\xi}&&
\end{aligned}
\]
\[
Q^n\qquad\left\{\begin{aligned}
&\xi_i\in H^{2i}(Q^n,\unsure{\mu|M}(i))\qquad\deleted{0\leq i\ \unsure{\text{pair}}\ \leq2n}\\
&\deleted{\lambda_{\unsure{2\nu|\nu}}\in H^{2\nu}(Q^{2\nu})}\quad\lambda_{2\nu}\in H^{2\nu}(Q^{2\nu})
\end{aligned}\right.
\]
\[
\deleted{Q}\qquad \xi_{\unsure{i|0}}
\]
\diagram{\begin{tikzcd}[ampersand replacement=\&,column sep=small]
Q^{n-1}\arrow[r,phantom,"\subset" description] \& Q^n\\
\xi_{\unsure{i|1}} \& \xi_i\arrow[l]\\
\xi_i\arrow[r,decorate,decoration={snake,amplitude=.26mm,segment length=1.7mm,post length=1.4mm}] \& \xi_{i+2}
\end{tikzcd}}
\[
\begin{pmatrix}1,1\\1,-1\end{pmatrix}
\]
\diagram{$\begin{array}{c@{\;}c@{\;}c@{\;}c@{\;}c@{\;}c@{\;}c@{\;}c@{\;}c}
Q^0&\subset&Q^1&\subset&Q^2&\subset&Q^3&\subset&Q^4\\
\lambda_0^{\unsure{*|0|+}},\xi^0_{\unsure{i|0|+}}&&\xi^0&&\xi^0&&\xi^0&&\xi^0\\
 &&\frac12\xi&&(\lambda_2,\xi)&&\xi&&\xi\\
 &&&&\frac12\xi^2&&\frac12\xi^2&&(\lambda_{\deleted{\unsure{2}}4},\xi^2)\\
 &&&&&&\frac12\xi^3&&\frac12\xi^3\\
 &&&&&&&&\frac12\xi^4
\end{array}$}
\[
\deleted{Q^{\unsure{n=2\nu}}}:\quad \deleted{\xi}\ \cdots
\]
% slanted
\margin{d'indice 2 dans $H^{2\nu}(Q^{2\nu})$, il faut \unsure{rajouter|ajouter} $\frac12(\xi^\nu+\lambda_\nu)=\varphi_\nu$ pour avoir une base $(\xi^\nu,\varphi_\nu)$}
\[
\boxed{\begin{aligned}
Q^{n=2\nu}:\quad &\xi^0\ \cdots,\xi^{\nu-1},[\xi^\nu,\lambda_\nu],\tfrac12\xi^{\nu+1},\cdots,\tfrac12\xi^{2\nu}\\
Q^{n=2\nu+1}:\quad &\xi^0,\cdots,\xi^\nu,\tfrac12\xi^{\nu+1},\cdots,\tfrac12\xi^{2\nu+1}
\end{aligned}}\qquad \lambda_\nu
\]
% archive sheet number: 23
